\chapter{Overview of the proof}\label{sec:overview}

In this section, we explain the proof of \cref{thm:main}, which
occupies \crefrange{sec:fnd}{sec:b7}. For each step, we state its main
estimate in the form in which it is proved, give the precise reference,
and explain briefly why the estimate holds and where it is used.
\Cref{sec:overview-map} contains a map of the proof, and
\cref{sec:overview-profiles} discusses three conventions that recur in
every part of it. Throughout, $\beta>1$ is fixed. We write
$X\lesssim Y$ if $X\le C_\beta Y$ for a constant $C_\beta$ that depends
only on $\beta$. Our estimates are uniform in the system size $N$, the interpolation
parameter $\rho$, the tilt $\lambda$, the minimizing profile, and the
positions and masses of its atoms. Much of the technical work goes into
this uniformity.

\section{The structure of the argument}\label{subsec:intro-outline}

We follow the interpolation scheme of~\cite{AL}; what changes is how
we estimate the self-overlap variance along the interpolation. The argument has four steps.

First, we reduce \cref{thm:main} to a bound on a fractional-moment
error $e_{N,\beta}(\lambda)$, which compares a normalized logarithmic
moment of order $\lambda/\beta$ of the partition function with a Parisi
problem constrained to have mass at least $\lambda/\beta$ at zero. By
\cref{prop:fnd-mean-bias}, the mean bias divided by $\beta$ is at most
$2e_{N,\beta}(\lambda)$ plus a term of order $\lambda^5$, and it
suffices to prove that $e_{N,\beta}(N^{-7/50})\lesssim N^{-7/10}$
(\cref{sec:overview-reduction}).

Second, we write $4e_{N,\beta}(\lambda)$ as the integral over
$\rho\in(0,1]$ of a defect $\calD_N(\rho)=\beta(A+2J+V)$ along an
interpolation, evaluated at a minimizing profile of a variational
problem in dimension $N$ (\cref{sec:overview-interpolation}). The terms
$A$ and $J$ are at most of order $\log^2(2+N)/(N\rho\lambda)$
(\cref{sec:overview-thermal}), and the difficulty is the self-overlap
variance $V$.

Third, we bound $V$ in two regimes. Near $\rho=0$, we combine estimates
on quantile intervals of the minimizing profile with a localization of
its lowest quantiles (\cref{sec:overview-shells}). For
$\rho$ above a cutoff $\rho_0$, we bound $V$ by a one-site cavity argument
(\cref{sec:overview-cavity}), which contracts if a scalar response
operator at the minimizing profile is at most $1+\rho/4$. We prove this
bound by showing that the operator is at most the identity at the Parisi
measure (\cref{sec:overview-stability}) and that the minimizing profile
is close to a scalar trial profile (\cref{sec:overview-localization}).
The two regimes give the following bounds. For almost every
$\rho\in[\rho_0,1]$, we have
\begin{equation}\label{eq:ov-stable}
 \calD_N(\rho)\lesssim\frac{\log^2(2+N)}{N\lambda\rho^2}
\end{equation}
by \cref{cor:c7-defect}, and for almost every $\rho\in(0,1]$, we have
\begin{equation}\label{eq:ov-small}
 \calD_N(\rho)\lesssim\frac{\log^2(2+N)}{N\rho\lambda}+(N\rho)^{-2/3}
 +(N\rho)^{-2/3}\lambda^{-1/3}\omega(\rho,\eta)^{2/3}
\end{equation}
by \cref{prop:sm-small}. Here $\eta=e_{N,\beta}(\lambda)+L_s\lambda^5$
is the trial budget of \cref{lem:fnd-trial}, with $L_s\ge1$ depending
only on $\beta$, and $\omega(\rho,\eta)=\min\{1,(\eta/\rho)^{1/6}\}$.

Fourth, the cutoff is $\rho_0=\min\{1,K\eta^{1/7}\}$, with $K$
depending only on $\beta$. Note that $\rho_0$ is defined through the
trial budget, which contains the unknown error. Integrating
\eqref{eq:ov-stable} over $[\rho_0,1]$ and \eqref{eq:ov-small} over
$(0,\rho_0]$, we obtain an implicit inequality whose leading term is
$N^{-2/3}\eta^{1/21}$. At $\lambda=N^{-7/50}$, this inequality
implies that $e_{N,\beta}(\lambda)\lesssim N^{-7/10}$
(\cref{sec:overview-closure}).

We take the first two steps from~\cite{AL}, where the same scheme
gives the rate $N^{-2/3}$~\cite[Theorem~1.4]{AL}, and our estimates
near $\rho=0$ refine those of~\cite[Sections~3 and~7]{AL}. We use the
cavity argument, and the stability and isolation estimates that it
requires, only above the cutoff.

\section{Reduction to the fractional-moment error}
\label{sec:overview-reduction}

As in~\cite[Section~1.1]{AL}, we complete the Hamiltonian by an
independent diagonal Gaussian term $D_N$, so that the completed
Hamiltonian has covariance exactly $NR^2/2$, where $R$ is the overlap
(\cref{sec:fnd-completion}). For $0<\lambda<\beta$, the
fractional-moment error is
\[
 e_{N,\beta}(\lambda)=E_\beta(\lambda)-\psi(\lambda),\qquad
 \psi(\lambda)=\frac1{N\lambda}\log\E_{\rm ref}
 e^{\lambda(F_{N,\beta}/\beta+D_N)},
\]
where $E_\beta(\lambda)=\beta^{-1}P_\beta(\lambda/\beta)$, and
$P_\beta(s)$ is the minimum of the Parisi functional over probability
measures $\mu$ on $[0,1]$ with $\mu(\{0\})\ge s$
\eqref{eq:fnd-error-definition}. Talagrand showed that the normalized
moments converge to this constrained
value~\cite{Talagrand2007FractionalMoments}; we need a rate.

The reduction is the midpoint argument of~\cite[Theorem~7.8,
Step~1]{AL}. Since $F_{N,\beta}/\beta+D_N$ is a convex Lipschitz function
of the Gaussian variables $(g_{ij},z_i)$, Chen's convexity
theorem~\cite[Theorem~1]{ChenConvexity} implies that $\psi$ extends to
a convex function on $[0,\infty)$ with
$\psi(0)=p_{N,\beta}/\beta$~\cite[Lemma~5.2]{AL}. Using
$\psi(0)\ge2\psi(\lambda)-\psi(2\lambda)$ and
$e_{N,\beta}(2\lambda)\ge0$, we find that
\[
 \frac{P_\beta-p_{N,\beta}}\beta
 \le E_\beta(0)-2E_\beta(\lambda)+E_\beta(2\lambda)
 +2e_{N,\beta}(\lambda).
\]
The second difference of $E_\beta$ is controlled by the cost of the
prescribed mass at zero. By~\cite[Lemma~7.2]{AL}, the Parisi profile
$\gamma_\beta=\beta\alpha_\beta$ grows linearly near $0$. As a consequence,
the scalar trial profile $\gamma^{\rm tr}_\lambda=\max\{\lambda,\gamma_\beta\}$
differs from it only on an interval of length of order $\lambda$, on
which the first variation vanishes at $\gamma_\beta$. The cost of the
trial profile is at most $L_s\lambda^5$
\eqref{eq:fnd-scalar-trial-cost}. Combining these facts, we arrive at
the bound of \cref{prop:fnd-mean-bias},
\[
 0\le\frac{P_\beta-p_{N,\beta}}\beta\le32L_s\lambda^5
 +2e_{N,\beta}(\lambda)\qquad(0<\lambda\le\lambda_0),
\]
with $\lambda_0>0$ depending only on $\beta$. We take
$\lambda=N^{-7/50}$ for $N$ so large that $\lambda\le\lambda_0$, and
then $\lambda^5=N^{-7/10}$. The tilt must be positive, because our
bounds on the defect degenerate as $\lambda\downarrow0$; we fix its exponent
in \cref{sec:overview-closure} by matching $\lambda^5$ with the
final bound on the error.

\section{The interpolation and the defect}
\label{sec:overview-interpolation}

The interpolation of~\cite[Section~2]{AL} is a finite-dimensional,
quantitative version of Guerra's interpolation~\cite{Guerra2003}, in the
form that Talagrand used for positive fractional
moments~\cite{Talagrand2007FractionalMoments}. For $0<\rho\le1$, let
$t=1-\rho$, $c=\beta^2\rho$, and $M=Nc$. An admissible profile is a
deterministic probability measure $\mu=s\delta_0+\nu$ on $[0,1]$, where
$s=\lambda/\beta$ and $\nu\ge0$, with distribution function
$a(q)=\mu([0,q])$. We solve the Parisi equation in $\R^N$,
\[
 \partial_qU=-\frac c2\bigl(\Delta U+a(q)|\nabla U|^2\bigr),\qquad
 U(1,x;W)=\log\sum_{\sigma\in\{-1,1\}^N}
 e^{\beta\sqrt t\,\sigma^TW\sigma/2+x\cdot\sigma},
\]
and we minimize the functional $\mathcal F_{t,s}(\mu)$ of
\eqref{eq:fnd-functional} over admissible profiles. At $t=0$, the
equation separates into $N$ scalar Parisi equations, and the minimum is
$P_\beta(s)$; at $t=1$, the minimum is the completed fractional moment
\eqref{eq:fnd-endpoints}. The fractional-moment error is
$\beta^{-1}$ times the decrease of the minimum along the interpolation.

\paragraph{The tilted law.}
Each profile defines a joint law of the couplings, an auxiliary field
$X_q\in\R^N$, $q\in[0,1]$, and a terminal spin $\sigma$. In this law, $W$ is
tilted by $e^{sU(0,0;W)}$, the field is a diffusion with drift
$ca(q)\nabla U(q,X_q;W)$ and noise coefficient $\sqrt c$, and $\sigma$ has
the Gibbs law in the external field $X_1$ (\cref{subsec:fnd-law}). We write $\calG_q$ for the
information revealed up to level $q$, $m_q$ and $C_q$ for the
conditional mean and covariance of $\sigma$ given $\calG_q$,
$H_q=\nabla^2U(q,X_q;W)$, and $S_q=|m_q|^2/N$. By
\cref{prop:fnd-contact}, $m_q$ is a martingale with
$\dd m_q=\sqrt c\,H_q\dd B_q$, and $a(q)C_q\preceq H_q\preceq C_q$. If
$\mu$ is a minimizer, every point $q$ of $\supp\nu$, called a contact
level, satisfies $\E S_q=q$ and $\E\Tr H_q^2\le N/c$, which are the
counterparts of the optimality conditions for the Parisi measure.

\paragraph{The defect.}
Let $R_q$ be the overlap of two replicas that are conditionally
independent given $\calG_q$. By \cref{prop:fnd-interpolation}, the
derivative of $\mathcal F_{t,s}(\mu)$ in $t$ is
$-(\beta^2/4)\int\E(R_q-q)^2\,\nu(\dd q)$ for every profile. It is a
square because the completed Hamiltonian has covariance $NR^2/2$, and
only $\nu$ appears because the outer moment of order $s$ cancels the
prescribed mass $s\delta_0$. At a contact level,
$R_q-S_q$ splits into three conditionally orthogonal pieces, which
gives $\E(R_q-q)^2=a^{\rm th}_q+2j_q+\Var(S_q)$ (\cref{lem:fnd-decomp}).
Integrating against $\nu$, we obtain \cref{cor:fnd-defect}, which
states that at almost every $\rho$ and at every minimizer,
\[
 \calD_N(\rho)=-\frac4\beta f_{N,s}'(1-\rho)=\beta(A+2J+V),
 \qquad 4e_{N,\beta}(\lambda)=\int_0^1\calD_N(\rho)\dd\rho,
\]
where $f_{N,s}(t)$ is the minimum of $\mathcal F_{t,s}$ and
\[
 A=\int\frac{\E\Tr C_q^2}{N^2}\,\nu(\dd q),\qquad
 J=\int\frac{\E[m_q^TC_qm_q]}{N^2}\,\nu(\dd q),\qquad
 V=\int\Var(S_q)\,\nu(\dd q).
\]
The terms $A$ and $J$ measure conditional fluctuations of the terminal
spin given $\calG_q$, while $V$ measures fluctuations of the information
that the field has revealed by level $q$. Each term may depend on the
minimizer, but their combination does not. We bound all three at an
arbitrary minimizer.

\paragraph{The posterior variance bound.}
Let $L_q=(\sigma\cdot m_q-|m_q|^2/2)/N$, whose conditional mean given
$\calG_q$ is $S_q/2$. By \cref{prop:fnd-variance}, we have, at a
minimizer and for every deterministic $k\in L^2(\nu)$,
\[
 \Var\Bigl(\int k(q)L_q\,\nu(\dd q)\Bigr)
 \le\frac1M\int\frac{k(q)^2}{a(q)}\,\nu(\dd q),
 \qquad A\le\frac1{Ms}.
\]
With the terminal spin fixed, the Gaussian coordinates that generate
$W$ and the field have a uniformly log-concave law, and the bound
follows from the Brascamp--Lieb inequality~\cite{BrascampLieb1976}, in
the form of~\cite[Proposition~2.2]{AL}, and the contact identities.
Note that the factor $a(q)^{-1}\le s^{-1}$ is where the prescribed mass
enters. We integrate such factors through the quantile function $Q$ of
$\nu$ \eqref{eq:fnd-quantile}, which pushes Lebesgue measure on
$(0,1-s)$ forward to $\nu$ and satisfies $a(Q(v))\ge s+v$.

\paragraph{The trial budget.}
In the energy normalization, we set $\gamma=\beta a$ and
$F_\rho(\gamma)=\beta^{-1}\mathcal F_{1-\rho,s}(\mu)$. By
\cref{lem:fnd-trial}, we have, for every $\rho\in(0,1]$,
\[
 0\le F_\rho(\gamma^{\rm tr}_\lambda)-\min_\gamma F_\rho(\gamma)
 \le e_{N,\beta}(\lambda)+L_s\lambda^5=\eta .
\]
Indeed, $\mathcal F_{t,s}(\mu)$ and its minimum are nonincreasing in
$t$, so that $F_\rho(\gamma^{\rm tr}_\lambda)$ is at most the scalar
value of the trial profile, which is at most $E_\beta(0)+L_s\lambda^5$,
while $\min F_\rho\ge\psi(\lambda)=E_\beta(\lambda)-e_{N,\beta}(\lambda)$
and $E_\beta(\lambda)\ge E_\beta(0)$. In words, the scalar trial profile has value at most $\eta$ above the minimum at every $\rho$. This budget controls
how far a minimizer can be from the trial profile
(\cref{sec:lz,sec:lz-root}), and it fixes the cutoff $\rho_0$
(\cref{sec:b7}). Since it contains the unknown error, the final inequality
is implicit.

\section{The thermal terms}\label{sec:overview-thermal}

At a contact level, by the covariance order, we have
$\Tr H_q^2\ge a(q)\Tr(H_qC_q)\ge a(q)^2\Tr C_q^2$, and the contact bound
on $\E\Tr H_q^2$ implies that $\E\Tr C_q^2/N^2\le(Ma(q)^2)^{-1}$. In
the quantile variable, we have
\[
 A\le\frac1M\int_0^{1-s}\frac{\dd v}{a(Q(v))^2}
 \le\frac1M\int_0^{1-s}\frac{\dd v}{(s+v)^2}\le\frac1{Ms}.
\]

The mean--covariance term needs more work, because
$j_q=N^{-2}\E[m_q^TC_qm_q]$ involves the conditional covariance in the
direction $m_q$, which is revealed at level $q$ itself. For a
direction $b$ that is measurable with respect to $\calG_{q'}$ for some
$q'<q$, we show by a scalar Riccati inequality that $\E[b^TC_qb]$ is at most
$\E|b|^2$ divided by the accumulated clock $c\int_{q'}^qa(u)^2\dd u$
(\cref{lem:th-ric}). However, this bound does not apply to $b=m_q$. To
overcome this obstacle, we write $m_{Q(v)}$ as a weighted average of
the earlier means $m_{Q(w)}$, $w<v$, plus increments. First-order
optimality implies that $\E S_r$ is $2$-Lipschitz toward every point
$q$ of $\supp\nu$, which bounds $\E|m_q-m_r|^2$ by $2N|q-r|$
(\cref{lem:th-gamma}). We control the increments by combining this
bound with \cref{lem:th-ric} in a dyadic decomposition (\cref{lem:th-res}).
The conditional covariance in the direction of the average is
controlled by the posterior variance bound. With a quadratic weight
$(s+w)^2$, both contributions are at most of order $(s+v)^{-2}$. We obtain
\[
 j_{Q(v)}\le\frac{341(1+\log_+M)^2}{M(s+v)^2},\qquad
 J\le\frac{341(1+\log_+M)^2}{Ms}
\]
(\cref{prop:th-j}).
Since $Ms=\beta N\rho\lambda$ and $M\le\beta^2N$, we conclude that
\[
 \beta(A+2J)\le\frac{1+682(1+\log_+M)^2}{N\rho\lambda}
 \lesssim\frac{\log^2(2+N)}{N\rho\lambda}.
\]
Integrated over $\rho\in[1/N,1]$, this is at most of order
$\log^3(2+N)/(N\lambda)$, which equals $N^{-43/50}\log^3(2+N)$ at
$\lambda=N^{-7/50}$ and is smaller than $N^{-7/10}$ by a power of $N$.

\section{The self-overlap variance for small \texorpdfstring{$\rho$}{rho}}
\label{sec:overview-shells}

The posterior variance bound controls averages of $L_q$ over levels.
However, $V$ involves $S_q$ at single levels, which may carry little
mass of $\nu$. Following~\cite[Lemma~3.3]{AL}, we compare $S_q$ with
averages of $L_p$ over windows of quantile levels next to $q$
(\cref{sec:wi}).

\paragraph{Two comparisons.}
For $p\ge q$, the conditional Jensen inequality for the martingale $m$
implies that $S_q\le2\E[L_p\mid\calG_q]$, and for $p\le q$, we have
$\E[L_p\mid\calG_q]=S_q/2-|m_q-m_p|^2/(2N)$. Averaging over a window
of quantile levels after or before $q$, with mean level $\bar p$, we
show in \cref{lem:wi-two} that, at a minimizer, $S_q-q$ is at most
(for a window after $q$) or at least (for a window before $q$) the sum
of $\bar p-q$ and twice the conditional expectation of a centered
average of $L_p$. If the window has quantile length $r$ and $a\ge u$ on
it, this conditional expectation has second moment at most
$(Mur)^{-1}$. Note that conditioning on $\calG_q$ removes the centered part
$N^{-1}(\sigma-m_q)\cdot m_q$ of $L_q$, which in~\cite[Lemma~3.3]{AL}
produces additional residual terms; we bound its variance $j_q$
separately (\cref{sec:overview-thermal}).

\paragraph{Variance on a quantile interval.}
Let $I\subset(0,1-s)$ be an interval of length $h$ such that
$a(Q(v))\ge u$ for $v\in I$, let $\ell$ be the diameter of $Q(I)$, and
set $\epsilon=(Mu)^{-1}$. For a threshold $\tau$, we take windows of
radius of order $\epsilon/\tau$, whose random parts have second moment
less than $\tau$. A point whose two windows lie in $I$ and for which
$\Var(S_{Q(v)})>\tau$ must then see an
increment of $Q$ at least of order $\tau^{1/2}$ across its window, and since $Q$
is nondecreasing with total increase $\ell$ on $I$, such points form a
set of measure at most of order $\epsilon\ell\tau^{-3/2}$. The remaining
points lie within distance of order $\epsilon/\tau$ of the ends of $I$. Integrating
$\min\{h,\epsilon\ell\tau^{-3/2}\}$ over $\tau$, we find the main term
in the bound of \cref{prop:wi-vfs},
\begin{equation}\label{eq:ov-vfs}
 \int_I\Var(S_{Q(v)})\dd v
 \le128\Bigl[\min\bigl\{h,\epsilon\bigl(1+\log_+(h/\epsilon)\bigr)\bigr\}
 +(\epsilon\ell)^{2/3}h^{1/3}\Bigr].
\end{equation}
The first term is at most $\epsilon$ up to a logarithm, and the
interval may contain part of the quantile interval of an atom.

\paragraph{Dyadic shells and root localization.}
In the proof of \cref{prop:sm-small}, we cover $(0,1-s)$ by dyadic
shells $I_u=\{v:u\le s+v<2u\}$, $u=2^js<1$, on which $a(Q(v))\ge u$,
$h\le u$, and $\epsilon=(Mu)^{-1}$. If we only use $\ell\le1$, the second
term in \eqref{eq:ov-vfs} is at most $M^{-2/3}u^{-1/3}$, and summing over
shells, with $\beta M^{-2/3}s^{-1/3}=(N\rho)^{-2/3}\lambda^{-1/3}$, we
would find only
\[
 \beta V\lesssim(N\rho)^{-2/3}\lambda^{-1/3}
 +\frac{\log^2(2+N)}{N\rho\lambda}.
\]
The bound $\beta V\le C(N\rho)^{-2/3}\lambda^{-1/3}$ that follows
from~\cite[(4.2)]{AL} has the same form. At $\lambda=N^{-7/50}$, the
factor $\lambda^{-1/3}=N^{7/150}$ would leave a positive power of $N$
after the normalization of \cref{sec:overview-closure}. To remove it, we
use the root localization of \cref{cor:lz-omega}, which states that at
every minimizer and for every trial budget $\eta$,
\[
 Q(v)\le C_Q\bigl(s+v+\omega(\rho,\eta)\bigr),\qquad
 \omega(\rho,\eta)=\min\{1,(\eta/\rho)^{1/6}\},
\]
with $C_Q$ depending only on $\beta$. Then
$\ell(I_u)\le2C_Q(u+\omega)$, the second term in \eqref{eq:ov-vfs} is
at most of order $M^{-2/3}(u^{1/3}+u^{-1/3}\omega^{2/3})$,
and its sum over shells is at most of order $M^{-2/3}(1+s^{-1/3}\omega^{2/3})$.
Together with the thermal terms, this proves \eqref{eq:ov-small}, that
is, \eqref{eq:sm-pointwise}. Its leading term $(N\rho)^{-2/3}$ does not
depend on $\lambda$, and the correction is integrable at $\rho=0$,
because $\omega^{2/3}\le(\eta/\rho)^{1/9}$. Below $\rho=1/N$, we rely
instead on the bound $\calD_N(\rho)\le C(N\rho)^{-1/2}$
of~\cite[Lemma~4.1]{AL}, recalled in \eqref{eq:fnd-moving}. Integrating,
we arrive at \eqref{eq:sm-integrated}, which states that for
$0<\rho_0\le1$,
\begin{equation}\label{eq:ov-small-integrated}
 \int_0^{\rho_0}\calD_N(\rho)\dd\rho
 \lesssim N^{-1}+\frac{\log^3(2+N)}{N\lambda}
 +N^{-2/3}\rho_0^{1/3}
 +N^{-2/3}\lambda^{-1/3}\eta^{1/9}\rho_0^{2/9}.
\end{equation}

\paragraph{\texorpdfstring{The barrier at $2/3$.}{The barrier at 2/3.}}
For $\rho_0=1$, the third term in \eqref{eq:ov-small-integrated} is of
order $N^{-2/3}$. This is the mechanism of~\cite[Proposition~7.7]{AL},
in which quantile estimates~\cite[Lemma~3.3]{AL} and a localization of
low quantiles~\cite[Lemma~7.6]{AL} lead to
$e_{N,\beta}(\lambda)\lesssim N^{-2/3}+N^{-3/4}\lambda^{-1/2}
+\log(2+N)/(N\lambda)$~\cite[(7.43)]{AL}, and the tilt
$\lambda=N^{-2/15}$ gives the rate $N^{-2/3}$
of~\cite[Theorem~1.4]{AL}. When $\rho$ is of order one, however,
$(N\rho)^{-2/3}$ is much larger than the cavity bound
\eqref{eq:ov-stable}, as we discuss heuristically in
\cref{subsec:intro-physics}. We therefore use
\eqref{eq:ov-small-integrated} only below the cutoff, where it
contributes $N^{-2/3}\rho_0^{1/3}$.

\section{The cavity contraction}\label{sec:overview-cavity}

Let $\mu$ be a minimizer at $(\rho,s)$ with distribution function $a$,
let $d=\beta^2(1-t/N)$, and let $J^{\scal}_{a,d}$ be the scalar response
operator of \cref{sec:overview-stability}, which acts on a subspace
$\mathcal H_\nu$ of $L^2(\nu)$. In \cref{prop:c7-closure}, we show that
if
\begin{equation}\label{eq:ov-stability}
 J^{\scal}_{a,d}\preceq(1+\eps_N)I\quad\text{on }\mathcal H_\nu,
 \qquad 0\le\eps_N\le\frac\rho4,
\end{equation}
then
\[
 \frac{3\rho}4V\le C_{\rm cav}\bigl[N^{-1}+A+J+H_3\bigr],\qquad
 H_3=\int\E|S_q-q|^3\,\nu(\dd q),
\]
with $C_{\rm cav}$ depending only on $\beta$. The argument occupies
\cref{part:cavity}, and \cref{sec:c7-mechanism} explains how the
inequality closes before the proofs begin.

\paragraph{Exchangeability.}
Let $n=N-1$, let $m_N(p)$ be the conditional magnetization of the last
spin given $\calG_p$, and let $A^-_p$ be the self-overlap of the
conditional mean of the other $n$ spins, which we call the rest
self-overlap. Then $NS_p=m_N(p)^2+nA^-_p$. Since the tilted law is
invariant under permutations of the sites, we have
$\Cov(S_p,m_N(p)^2)=\Var(S_p)$ \eqref{eq:c7-exchange}, and it follows
that
\[
 \Var(S_p)=\frac1N\Var\bigl(m_N(p)^2\bigr)
 +\frac nN\Cov\bigl(m_N(p)^2,A^-_p\bigr).
\]
Up to an error of order $N^{-1}$, the variance $V$ is the integrated
covariance between the squared magnetization of one spin and the
self-overlap of the others. Exchangeability is used in the same way in
the cavity method of~\cite[Section~1.6]{TalagrandVolI2011}.

\paragraph{The cavity interpolation.}
For $u\in[0,1]$, we give each coupling between the last spin and the
others variance $\alpha u/n$, and the last coordinate of the field
variance coefficient $b(u)=d-\alpha u$, where $\alpha=\beta^2tn/N$ and
$d=c+\alpha$ \eqref{eq:cavityparameters}. At $u=1$, we recover the
tilted law of the full system. At $u=0$, the last spin is independent
of the others: the law is a product of the tilted law of an $n$-spin
system, with the same profile, and the scalar law of \cref{sec:sc} with
coefficient $d$. In particular,
$\mathcal C_{pp}(u)=\Cov_u(m_N(p)^2,A^-_p)$ vanishes at $u=0$. By Taylor's formula, we
have
\[
 \mathcal C_{pp}(1)=\mathcal C_{pp}'(0)
 +\int_0^1(1-u)\,\mathcal C_{pp}''(u)\dd u .
\]
Note that the $n$-spin system at $u=0$ carries the profile of the full
system, which need not minimize its own functional.

\paragraph{Tree laws and the remainder.}
Derivatives in $u$ produce products of conditional expectations at
several levels. We represent these products by laws of finitely many
descendants on marked trees, in which deleting a leaf does not change
the law of the others (\cref{subsec:c7-tree}). Normalized replica trees of this
kind were developed in~\cite[Section~2]{OpenAIFluctuationScale}. By
\cref{lem:streaming}, the derivative in $u$ of a tree moment is
$\alpha/2$ times a combination of moments with one more overlap defect
$D_{\mathsf{ab}}=\epsilon_{\mathsf a}\epsilon_{\mathsf b}(R^-_{\mathsf{ab}}-q_{\mathsf{ab}})$,
plus a term from the normalization of the tilt. The row contributes the
rest overlap $R^-_{\mathsf{ab}}$ of two leaves, and the decrease of
$b(u)$ subtracts the level $q_{\mathsf{ab}}$ at which they split. It
follows that nonnegative tree moments at different values of $u$ agree
up to a factor $K_{\beta,k}$ that depends only on $\beta$ and the
number $k$ of leaves (\cref{cor:comparison}), and that every term of
$\mathcal C_{pp}''$ contains three defects. By
\cref{lem:cav-threeproducts,lem:remainder,lem:c7-thermal}, the Taylor
remainder is at most $C_\beta\sup_u\mathcal M_3(u)$, where
$\mathcal M_3(u)=\int\E_{2,q,u}|R^-_{12}-q|^3\,\nu(\dd q)$, and this
supremum is at most $K_{\beta,2}[64(A+J)+16H_3+32N^{-3}]$.

\paragraph{The linear term.}
Along the interpolation, the fields acting on the last spins of two
leaves have covariance
$dq_{\mathsf{ab}}+\alpha u(R^-_{\mathsf{ab}}-q_{\mathsf{ab}})$
(\cref{sec:c7-mechanism}). In the scalar law, if we replace the covariance $d(q\wedge r)$ of the
field by $dT(q\wedge r)$ with $T(q)=q+\theta\varphi(q)$, the mean square
magnetization at level $p$ changes at the rate $(J_{a,d}\varphi)(p)$
(\cref{lem:sc-response}). Heuristically, the row acts on the last spin
as a random transport of levels with displacement
$(\alpha u/d)(A^-_q-q)$, and we expect the main term of
$\mathcal C_{pp}'(0)$ to be $(\alpha/d)(J^{\scal}_{a,d}\Sigma^-)(p,p)$,
where $\Sigma^-(p,q)=\Cov^-(A^-_p,A^-_q)$ is the covariance of the rest
self-overlap at $u=0$. In \cref{lem:linear}, we prove the exact
identity
\[
 \int\mathcal C_{pp}'(0)\,\nu(\dd p)
 =\alpha\Bigl[\frac1d\Tr_\nu(J^{\scal}_{a,d}\Sigma^-)-\frac12\mathcal P
 +\int T_p\,\nu(\dd p)\Bigr],\qquad\mathcal P\ge0 .
\]
The term $-\mathcal P/2$ comes from the difference between conditional
and unconditional covariances of the rest self-overlap; we discard
it. The errors $T_p$ raise an obstacle: they contain covariances of
centered variables of the rest system with $A^-_p$, and if we applied the
Cauchy--Schwarz inequality at a single level, the bound would involve
$\Var^-(A^-_p)$, whose integral is what we want to control. To overcome this obstacle, we first integrate in $p$ against
the deterministic kernels of the scalar law, so that the covariance
order and the bound \eqref{eq:cl-BL} on the variance of $\nu$-averages
of $A^-$ apply. After a comparison with the full system through
\cref{cor:comparison}, we find in \cref{lem:cl-response} that
\[
 \alpha\Bigl|\int T_p\,\nu(\dd p)\Bigr|
 \le192\beta^2K_{\beta,3}\bigl[A+J+N^{-2}\bigr].
\]

\paragraph{\texorpdfstring{The factor $1-\rho$.}{The factor 1-rho.}}
Combining these identities, we obtain the trace identity of
\cref{lem:cav-trace},
\[
 V=\frac{V^{\rm site}}N+\frac nN\Bigl[
 \frac\alpha d\Tr_\nu(J^{\scal}_{a,d}\Sigma^-)
 -\frac\alpha2\mathcal P+\alpha\int T_p\,\nu(\dd p)
 +\mathcal R_{\rm cav}\Bigr],
\]
where $0\le V^{\rm site}\le1/4$ and $\mathcal R_{\rm cav}$ is the Taylor
remainder. The operator $\Sigma^-$ is positive, with range in
$\mathcal H_\nu$ and trace $V^-=\int\Var^-(A^-_p)\,\nu(\dd p)$. Under
\eqref{eq:ov-stability}, the trace $\Tr_\nu(J^{\scal}_{a,d}\Sigma^-)$ is
therefore at most $(1+\eps_N)V^-$. The contraction comes from the ratio
\[
 \frac\alpha d=t\,\frac{N-1}{N-t}\le t=1-\rho
\]
\eqref{eq:c7-alpha-d}. Indeed, the field acting on the last spin has
variance coefficient $d$, and at $u=1$, the part $c=\beta^2\rho$ of it
comes from the auxiliary field, whose covariance is deterministic; only
the part $\alpha$ that comes from the couplings feels the fluctuations
of the rest overlap. We conclude that
$V\le(1-\rho)(1+\eps_N)V^-+{}$errors.

\paragraph{Comparison with coefficient one.}
It remains to compare $V^-$ with $V$, and the comparison must have
coefficient one. If we compared the endpoints through
\cref{cor:comparison}, we would obtain only $V^-\le KV+{}$errors with
some $K>1$, and the inequality $V\le(1-\rho)K(1+\eps_N)V+{}$errors bounds $V$
only if $(1-\rho)K(1+\eps_N)<1$, which fails for $\rho\le1-1/K$. In
\cref{lem:cl-star}, we prove instead that
\[
 V^-\le V+C_\beta\sup_{0\le u\le1}\mathcal M_3(u)+\frac1{4n^2}.
\]
For each $p$, the second moment $\E_u(A^-_p-p)^2$ is the mean of a
product of two overlap defects on four leaves that split at $p$. Its
derivative in $u$ contains three defects, and after integration against
$\nu$, the values at $u=0$ and $u=1$ differ by at most
$C_\beta\sup_u\mathcal M_3(u)$. At $u=0$, it bounds $\Var^-(A^-_p)$, even
though we do not know that $\E^-A^-_p=p$. At $u=1$, exchangeability and
the contact identity $\E S_p=p$ show that it is at most
$\Var(S_p)+1/(4n^2)$.

\paragraph{Closing.}
As $(1-\rho)(1+\rho/4)\le1-3\rho/4$, these bounds imply that
$V\le(1-3\rho/4)V+C_\beta[N^{-1}+A+J+H_3]$, which is
\cref{prop:c7-closure}. By \cref{cor:wi-cub}, if $M\ge1$, then
\[
 H_3\le zV+\frac{425}{Ms}\log(2M)\log\frac2z\qquad(0<z\le1/4).
\]
We take $z$ proportional to $\rho$ and absorb $zV$ in the margin
$3\rho/4$. Inserting the bounds on $A$ and $J$ and using
$Ms=\beta N\rho\lambda$, we arrive at \cref{cor:c7-defect}, which states
that at every minimizer that satisfies \eqref{eq:ov-stability}, we have
\[
 \beta(A+2J+V)\le C_\beta\frac{\log^2(2+N)}{N\lambda\rho^2}
 \qquad(1/N\le\rho\le1).
\]
One factor $\rho^{-1}$ comes from the margin, and the other from $Ms$.

The framework of differentiation, comparison, and contraction on one
site is that of~\cite[Section~1.6]{TalagrandVolI2011}. A one-site
expansion with an independent scalar endpoint also appears
in~\cite[Lemma~2.2]{OpenAILimitingLaw}, where the resulting linear
equation is inverted under a gap condition on the scalar
response~\cite[Proposition~5.1]{OpenAILimitingLaw}; we compare the two
arguments in \cref{subsec:intro-openai}. We use no gap condition. The
multiplier that dominates our scalar response equals one on the support
of the Parisi measure, and the contraction comes from the ratio
$\alpha/d$.

\section{Stability of the scalar response}\label{sec:overview-stability}

At the independent endpoint, the last spin follows the scalar law with
the profile $a_N$ of the $N$-spin minimizer and the coefficient
$d_N=\beta^2(1-t/N)$, not with the Parisi measure. In \cref{sec:sc}, we
verify \eqref{eq:ov-stability} with an explicit $\eps_N$.

\paragraph{The response operator.}
For a profile $a$ and $d\in[\beta^2/2,\beta^2]$, let $u=u_{a,d}$ solve
the scalar Parisi equation with coefficient $d$ \eqref{eq:sc-pde}, and
let $Y$ be its diffusion. With $m_q=\partial_xu(q,Y_q)$ and
$h_q=\partial_{xx}u(q,Y_q)$, we define
\[
 \Gamma_{a,d}(q)=\E^0m_q^2,\qquad k_{a,d}(q)=d\,\E^0h_q^2,
\]
so that $\Gamma_{a,d}'=k_{a,d}$ \eqref{eq:sc-martingale}; we call
$k_{a,d}$ the scalar multiplier. By \cref{lem:sc-response}, the
operator $J_{a,d}$ of \eqref{eq:scalarJexplicit} is the derivative of
$\Gamma_{a,d}$ under transports of covariance levels: if we replace the
covariance $d(q\wedge r)$ of the field by $dT_\theta(q\wedge r)$, where
$T_\theta(q)=q+\theta\varphi(q)$ and $\varphi(0)=\varphi(1)=0$, the mean
square magnetization at level $p$ has derivative $(J_{a,d}\varphi)(p)$
at $\theta=0$. The scalar response operator $J^{\scal}_{a,d}$ is the
compression of $J_{a,d}$ to the space $\mathcal H_\nu$ of functions in
$L^2(\nu)$ that vanish at $0$ and $1$ when these are atoms of $\nu$. Its
kernel has no sign, and a direct estimate bounds its norm only by
$C_\beta$.

\paragraph{Domination by the multiplier.}
The following observation reduces the operator bound to a pointwise
bound. In \cref{lem:Jdom}, we show that, for every profile,
$J^{\scal}_{a,d}$ is self-adjoint and
\[
 J^{\scal}_{a,d}\preceq M_{k_{a,d}},\qquad
 (M_{k_{a,d}}v)(q)=k_{a,d}(q)v(q).
\]
To prove it, we follow the scalar Parisi functional $\mathcal P_d$ along
a transport of levels, $\Pi(\theta)=\mathcal P_d(a\circ T_\theta^{-1})$.
By \cref{lem:sc-response}, we have
$\Pi''(0)=(d/2)\langle\varphi,(I-J_{a,d})\varphi\rangle$
\eqref{eq:transporthess}; the first derivative is the first variation
under transports of levels due to
Talagrand~\cite{Talagrand2006ParisiMeasures}
(see~\cite[Section~3, (19)]{AuffingerChen2015Properties}). Since the functional is convex along segments of distribution
functions~\cite{AuffingerChen2015Uniqueness}, its increment along the
transport is at least its first variation at $a$, applied to
$a\circ T_\theta^{-1}-a$. We expand this lower bound to second order in
$\theta$, which produces the integral of $(1-k_{a,d})\varphi^2$ against
$\nu$. Comparing the two expansions, we find that
$\langle\varphi,J_{a,d}\varphi\rangle\le\langle\varphi,k_{a,d}\varphi\rangle$.
Note that $\Pi$ need not be convex, because
$\theta\mapsto a\circ T_\theta^{-1}$ is not affine; the term
$1-k_{a,d}$ accounts for this. It therefore suffices to show that
$\sup_qk_{a_N,d_N}(q)\le1+\eps_N$.

\paragraph{The multiplier at the Parisi measure.}
In \cref{lem:reference}, we prove that
\[
 k_{\alpha_\beta,\beta^2}(q)=1\quad(0\le q\le q_\beta),\qquad
 k_{\alpha_\beta,\beta^2}(q)\le e^{-\beta^2(q-q_\beta)/2}
 \quad(q_\beta\le q\le1).
\]
On the support, we rely on an identity of
Talagrand~\cite{Talagrand2006ParisiMeasures}, in the form
of~\cite[Theorem~5]{AuffingerChen2015Properties}
and~\cite[Lemma~2.6]{Lopatto}, which states that $\Gamma(q)=q$ on
$\supp\mu_\beta$. As the support is the interval
$[0,q_\beta]$~\cite[Theorem~1.1]{Lopatto}, we may differentiate this
identity along it, which gives $k=\Gamma'=1$. Above the support,
we have $\alpha_\beta=1$, and the solution $u$ is explicit. A
transformed density $z_q$ of $Y_q$ evolves there by the heat
equation~\cite[Lemma~2.4(ii)]{Lopatto}, and $z_{q_\beta}$ is even and
log-concave~\cite[Proposition~4.1]{Lopatto}. In particular,
$z_{q_\beta}$ is decreasing on $(0,\infty)$. The multiplier is an
explicit multiple of the integral of $z_q$ against $\sech^3$
\eqref{eq:sc-k-plateau}, and Gaussian smoothing of an even decreasing
density does not increase its integral against an even decreasing
function; this gives the exponential bound. We need it on all of
$[0,1]$, because a finite-dimensional minimizer may place mass above
$q_\beta$.

\paragraph{Marginality.}
At the Parisi measure, the multiplier equals one on the whole support,
and the domination gives no contraction there. The cavity
inequality contracts because of the ratio $\alpha/d\le1-\rho$, which
leaves a margin of order $\rho$ only. This is why we need $\eps_N\le\rho/4$ in
\eqref{eq:ov-stability}, and why stability holds only above a cutoff.
The equality $k=1$ on the support expresses, in scalar form, the
marginal stability that De~Dominicis and Kondor found in a truncated
replica model near the critical
temperature~\cite[pp.~606--608]{DeDominicisKondor1983}; see
\cref{subsec:intro-physics}.

\paragraph{Transfer to the minimizer.}
By \cref{lem:scalarcontinuity}, the multiplier depends Lipschitz
continuously on the profile and coefficient, in the sense that
\[
 \sup_{0\le q\le1}|k_{a,d}(q)-k_{\widetilde a,\widetilde d}(q)|
 \le C_\beta\bigl(\|a-\widetilde a\|_{L^1(0,1)}+|d-\widetilde d|\bigr).
\]
We have $|d_N-\beta^2|\le\beta^2/N$ and
$\|\gamma^{\rm tr}_\lambda-\gamma_\beta\|_{L^1(0,1)}\lesssim\lambda^2$
\eqref{eq:sc-floor}. Combining these facts with the isolation estimate
of \cref{sec:overview-localization}, we obtain \cref{prop:lz-stability},
which states that on the space $\mathcal H_{\nu_N}$ of the minimizer
itself,
\[
 J^{\scal}_{a_N,d_N}\preceq(1+\eps_N)I,\qquad
 \eps_N\le C_{\rm st}\bigl[\rho^{-3/4}\eta^{1/4}+\lambda^2+N^{-1}\bigr],
\]
with $C_{\rm st}$ depending only on $\beta$. We compare only the multipliers of
the two profiles, and the operator inequality comes from
\cref{lem:Jdom} at $a_N$.

\section{Localization of the minimizing profile}
\label{sec:overview-localization}

We need two bounds on the distance between a minimizer $\gamma_*$ of
$F_\rho$ and the scalar trial profile $\gamma^{\rm tr}$, in terms of a
trial budget $\eta$: a bound in $L^1(0,1)$ for the stability of the
scalar response (\cref{sec:lz}), and a localization of the lowest
quantiles for small $\rho$ (\cref{sec:lz-root}).

\paragraph{Convexity and response.}
From~\cite[Section~7.4 and Appendix~A]{AL}, we import the following
facts (\cref{prop:loc-inputs}). The functional $F_\rho$ is convex along
segments $\gamma_\theta=(1-\theta)\gamma_0+\theta\gamma_1$. Let
$f(\theta)=F_\rho(\gamma_\theta)$
and $\zeta_\theta(q)=\beta^2N^{-1}\E_{\gamma_\theta}\Tr H_q^2$, so that
$\rho\zeta_\theta$ is the derivative of the expected self-overlap. The
second variation controls how the self-overlap responds to
$\gamma_1-\gamma_0$:
\[
 \rho\left|\int_0^1(\gamma_1-\gamma_0)\zeta_\theta\varphi\dd q\right|
 \le\Bigl(\frac{f''(\theta)}\rho\Bigr)^{1/2}
 \bigl[Y^{1/2}+2(Y+Z)^{1/2}\bigr]
\]
\eqref{eq:loc-response}, for Lipschitz cutoffs $\varphi$, where $Y$ is a
weighted Dirichlet energy of $\varphi$ and $Z$ is the integral of
$\rho\varphi^2\zeta_\theta$ against $\dd\gamma_\theta$. Further,
$\zeta_\theta\ge k_\beta^2$ for a constant $k_\beta>0$
\eqref{eq:loc-in-curv}. The midpoint inequality \eqref{eq:loc-midpoint}
bounds the integral of $\min\{\theta,1-\theta\}f''(\theta)$ over
$(0,1)$ by $F_\rho(\gamma_0)+F_\rho(\gamma_1)-2\min F_\rho$. If both
endpoints are nearly minimal, the segment has little
curvature, while a difference of one sign and of definite size on some
region forces curvature (\cref{lem:loc-coercive}).

\paragraph{\texorpdfstring{Isolation in $L^1$.}{Isolation in L1.}}
A minimizer may have atoms, and $Z$ contains $\dd\gamma_\theta$ on the
support of the cutoff. To overcome this difficulty, we first spread the
free mass of $\gamma_*$ on a scale $h$ by a reflected regularization
(\cref{lem:loc-smoothing}). It creates a density at most $\beta/h$ at a
cost at most $\rho\beta h^2/12$; the cost is quadratic because the
first-variation potential is semiconcave, as in~\cite[Lemma~7.5]{AL}.
Along the segment from the regularized minimizer to $\gamma^{\rm tr}$,
we use cutoffs of one sign built from one-sided averages of
superlevel sets of the difference, away from the atom of
$\gamma^{\rm tr}$ at $q_\beta$; below $q_\beta$, the trial profile is
Lipschitz (\cref{lem:loc-lipschitz}). Using the layer-cake formula, we
find that $f''\gtrsim\rho^2\min\{\rho D^4,hD^2\}$, where
$D=\|\gamma_*-\gamma^{\rm tr}\|_{L^1(0,1)}$ \eqref{eq:loc-E2-second}.
For $h$ proportional to $\rho D^2$, we have $f''\gtrsim\rho^3D^4$, and
the smoothing cost is a small fraction of $\rho^3D^4$. The midpoint
inequality implies that $\rho^3D^4\lesssim\eta$, which is the bound of
\cref{prop:loc-isolation},
\[
 \|\gamma_*-\gamma^{\rm tr}\|_{L^1(0,1)}
 \le C\min\{1,\rho^{-3/4}\eta^{1/4}\},
\]
for every minimizer and every trial budget, with $C$ depending only on
$\beta$.

\paragraph{Root localization.}
In the isolation argument, the lower bound
$\rho\int\zeta_\theta\varphi\ge\rho k_\beta^2\int\varphi$ on the response
carries a factor $\rho$, which is too costly near $\rho=0$. To avoid it,
we use the local deficit certificate of \cref{lem:lz-certificate}. Suppose that
two nearly minimal profiles have bounded densities on an interval $I$ of
length $\ell$, that the density of the second is also bounded below,
and that the second exceeds the first by a definite amount on $I$. Then
the certificate bounds $\eta/\rho$ from below by an explicit expression
in $\ell$, the gap, and the density bounds. In
its proof, we move mass of the profile toward the middle of $I$ along a
bump. Since both endpoints are nearly minimal and $F_\rho$ is convex,
the response integrated against the bump must be at least of order
$\ell$, without the factor $\rho$. To localize
a quantile $Q_{\gamma_*}(v)$, we bound its truncation
$q=\min\{Q_{\gamma_*}(v),T_1\}$ at the level $T_1$ of
\cref{lem:lz-rootdensity}, which depends only on $\beta$; this suffices,
because $Q_{\gamma_*}(v)\le q/T_1$. If $q$ is large compared with
$\lambda+v$, we regularize the free mass of the minimizer below $q$ only. The result
is small on $(q/4,3q/4)$, where the trial profile has density at least
$k_0/2$ (\cref{lem:lz-rootdensity}) and exceeds it by an amount of order
$q$. Using the certificate, we find that
$\eta/\rho\gtrsim\min\{q^6,q^4/(m/h+C)\}$, where $m/h$ is the density
created by the regularization. With a width adapted to the moved mass,
we obtain \cref{prop:lz-quantile},
which states that
\[
 Q_{\gamma_*}(v)\le C\bigl[\lambda+v+\min\{1,(\eta/\rho)^{1/6}\}\bigr]
 \qquad(0<v<\beta-\lambda),
\]
together with its form \eqref{eq:lz-Q} in the probability
normalization. We follow the strategy of~\cite[Lemma~7.6]{AL}. Compared
with~\cite[(7.30)]{AL}, the bound depends on $\eta$ and $\rho$ only
through $\eta/\rho$, and the certificate is what gives this dependence.

\section{Closing the implicit inequality}\label{sec:overview-closure}

By \cref{prop:lz-stability}, the stability assumption
\eqref{eq:ov-stability} holds as soon as
$C_{\rm st}[\rho^{-3/4}\eta^{1/4}+\lambda^2+N^{-1}]\le\rho/4$. Recall
that $\eta\ge L_s\lambda^5\ge\lambda^5$, so that
$\eta^{1/7}\ge\lambda^{5/7}=N^{-1/10}$, which is at least $\lambda^2$ and
$N^{-1}$. The first term satisfies $\rho^{-3/4}\eta^{1/4}\lesssim\rho$
exactly when $\rho\gtrsim\eta^{1/7}$. With $K=12C_{\rm st}$, we set
\[
 \rho_0=\min\{1,K\eta^{1/7}\}
\]
\eqref{eq:b7-cutoff}. If $\rho_0<1$, then $\eps_N\le\rho/4$ for every
$\rho\in[\rho_0,1]$ and every minimizer, and \eqref{eq:ov-stable}
applies; if $\rho_0=1$, the stable interval is empty. Using
$\rho_0\ge\lambda^{5/7}$, we find that
\[
 \int_{\rho_0}^1\calD_N(\rho)\dd\rho
 \lesssim\frac{\log^2(2+N)}{N\lambda\rho_0}
 \lesssim\frac{\log^2(2+N)}{N\lambda^{12/7}}
 =N^{-19/25}\log^2(2+N).
\]
On $(0,\rho_0]$, we insert $\rho_0\le K\eta^{1/7}$ into
\eqref{eq:ov-small-integrated}, which gives
$\rho_0^{1/3}\lesssim\eta^{1/21}$ and
$\eta^{1/9}\rho_0^{2/9}\lesssim\eta^{1/7}$. Adding the two
contributions and recalling that
$4e_{N,\beta}(\lambda)=\int_0^1\calD_N(\rho)\dd\rho$, we obtain the
bound \eqref{eq:b7-normalized-input},
\begin{equation}\label{eq:ov-implicit}
 e_{N,\beta}(\lambda)\lesssim N^{-19/25}\log^3(2+N)
 +N^{-2/3}\eta^{1/21}+N^{-2/3}\lambda^{-1/3}\eta^{1/7}.
\end{equation}

But there is an apparent circularity here, since the cutoff $\rho_0$ is
defined through $\eta=e_{N,\beta}(\lambda)+L_s\lambda^5$, which contains
the quantity that we want to bound. The way out is that the error
enters \eqref{eq:ov-implicit} only through powers of $\eta$ smaller
than one. Let $r=N^{-7/10}=\lambda^5$ and $x=\eta/r$, which is finite
and at least $L_s\ge1$. Dividing \eqref{eq:ov-implicit} by $r$ and
adding $L_s$ to both sides, we find \eqref{eq:b7-normalized},
\[
 x\le L_s+C_\beta\bigl[N^{-3/50}\log^3(2+N)+x^{1/21}
 +N^{-1/50}x^{1/7}\bigr]\le C_\beta x^{1/7},
\]
which implies that $x\le C_\beta^{7/6}$ and
$e_{N,\beta}(\lambda)\le\eta=rx\lesssim N^{-7/10}$. We need no preliminary bound
on the error. The proof of~\cite[Proposition~7.7]{AL} closes
an inequality in which the unknown error enters the trial budget in the
same way. By \cref{prop:fnd-mean-bias}, we conclude that
$P_\beta-p_{N,\beta}\lesssim N^{-7/10}$ for all large $N$, and the
finitely many remaining values of $N$ only change the constant.

\paragraph{The exponent.}
If $\eta$ is of order $N^{-a}$, the leading term
$N^{-2/3}\eta^{1/21}$ in \eqref{eq:ov-implicit} is of the same order as
$\eta$ when
\[
 a=\frac23+\frac a{21},\qquad\text{that is,}\qquad a=\frac7{10},
\]
and the cost $\lambda^5$ is of the same order when $\lambda=N^{-7/50}$
(\cref{rem:b7-exponents}). \Cref{tab:overview-terms} lists the terms of
\eqref{eq:ov-implicit} after division by $N^{-7/10}$. Only the leading
term has no polynomial slack; the others carry the factors $N^{-1/50}$,
$N^{-3/50}$, $N^{-4/25}$, and $N^{-3/10}$, up to logarithms. After the
closure, we have $L_sN^{-7/10}\le\eta\lesssim N^{-7/10}$, and the
cutoff is then comparable to $N^{-1/10}$ for large $N$.
\Cref{tab:overview-exponents} collects the exponents of the proof and
the statements in which they arise.

\begin{table}[!tb]
\centering
\small
\begin{tabular}{@{}>{\raggedright\arraybackslash}p{0.27\textwidth}
 >{\raggedright\arraybackslash}p{0.42\textwidth}
 >{\raggedright\arraybackslash}p{0.23\textwidth}@{}}
\toprule
Term & Origin & After division by $N^{-7/10}$\\
\midrule
$N^{-1}$ & $\rho\le1/N$, by \eqref{eq:fnd-moving} & $N^{-3/10}$\\
$N^{-43/50}\log^3(2+N)$ & thermal and logarithmic terms on
 $[1/N,\rho_0]$, \eqref{eq:sm-pointwise} & $N^{-4/25}\log^3(2+N)$\\
$N^{-19/25}\log^2(2+N)$ & stable interval $[\rho_0,1]$,
 \cref{cor:c7-defect} & $N^{-3/50}\log^2(2+N)$\\
$N^{-2/3}\eta^{1/21}$ & quantile shells on $(0,\rho_0]$,
 \cref{prop:wi-vfs} & $x^{1/21}$\\
$N^{-2/3}\lambda^{-1/3}\eta^{1/7}$ & root localization correction,
 \cref{cor:lz-omega} & $N^{-1/50}x^{1/7}$\\
\bottomrule
\end{tabular}
\caption{The terms of the bound on $e_{N,\beta}(\lambda)$ at
$\lambda=N^{-7/50}$, with $x=\eta/N^{-7/10}$. The first three terms are
combined into one in \eqref{eq:ov-implicit}; see
\cref{subsec:b7-general}.}
\label{tab:overview-terms}
\end{table}

\begin{table}[!tb]
\centering
\small
\begin{tabular}{@{}>{\raggedright\arraybackslash}p{0.2\textwidth}
 >{\raggedright\arraybackslash}p{0.47\textwidth}
 >{\raggedright\arraybackslash}p{0.25\textwidth}@{}}
\toprule
Quantity & Where it comes from & Statement\\
\midrule
$\lambda^5$ & cost of the prescribed mass at zero, at the scalar end
 & \eqref{eq:fnd-scalar-trial-cost}, \cite[Lemma~7.2]{AL}\\
$(N\rho)^{-2/3}$ & variance on dyadic quantile shells, after root
 localization & \cref{prop:wi-vfs}, \eqref{eq:sm-pointwise}\\
$(\eta/\rho)^{1/6}$ & localization of the lowest quantiles
 & \cref{prop:lz-quantile}, \cref{cor:lz-omega}\\
$\rho^{-3/4}\eta^{1/4}$ & $L^1$ isolation of the minimizing profile
 & \cref{prop:loc-isolation}\\
$\rho/4$ & admissible error in the scalar response; it leaves the
 margin $3\rho/4$ & \eqref{eq:c7-stab}\\
$\eta^{1/7}$ & cutoff: $\rho^{-3/4}\eta^{1/4}\lesssim\rho$ when
 $\rho\gtrsim\eta^{1/7}$ & \eqref{eq:b7-cutoff}\\
$(N\lambda\rho^2)^{-1}$ & cavity bound: margin $3\rho/4$ and
 $Ms=\beta N\rho\lambda$ & \cref{cor:c7-defect}\\
$\rho_0^{1/3}$ & integral of $(N\rho)^{-2/3}$ over $(0,\rho_0]$
 & \eqref{eq:sm-integrated}\\
$\eta^{1/21}$ & $\rho_0^{1/3}$ with $\rho_0\le K\eta^{1/7}$
 & \eqref{eq:b7-normalized-input}\\
$a=7/10$ & balance $a=2/3+a/21$ & \cref{rem:b7-exponents}\\
$\lambda=N^{-7/50}$ & $\lambda^5=N^{-7/10}$ & \eqref{eq:b7-parameters}\\
\bottomrule
\end{tabular}
\caption{The exponents of the proof.}
\label{tab:overview-exponents}
\end{table}

\paragraph{What would improve the exponent.}
The exponent $1/21$ is the product of the exponent $1/3$ of $\rho_0$ in
the integrated shell estimate and the exponent $1/7$ of $\eta$ in the
cutoff, which comes from the isolation estimate. A larger exponent in
either place would improve the balance. A smaller cutoff, however,
increases the contribution $\log^2(2+N)/(N\lambda\rho_0)$ of the stable
interval, and such an improvement requires rechecking the other
terms of \cref{tab:overview-terms}. If the cavity bound
\eqref{eq:ov-stable} held down to $\rho$ of order $N^{-1/4}$, and if we
ignored logarithms and the tilt, the two regimes would give $N^{-3/4}$;
we discuss this heuristic in \cref{subsec:intro-physics}.

\section{Minimizers, arbitrary profiles, and atoms}
\label{sec:overview-profiles}

Three conventions recur throughout the proof.

\paragraph{Which statements use minimality.}
Estimates that use the contact identities \eqref{eq:in-contact} or the
trial budget hold at minimizers: the main results of
\cref{sec:wi,sec:th,sec:lz,sec:lz-root}, Step~2 of the proof of
\cref{lem:cl-star}, and everything built on them. The Gaussian tools of
\cref{sec:gauss}, the lemmas of \cref{sec:sc} on the scalar law of an
arbitrary profile, and the cavity calculus of
\cref{sec:c7,sec:c7-endpoint} hold for every deterministic profile (with
$\nu(\{1\})=0$ in \cref{sec:c7-endpoint}). The reference estimate
\cref{lem:reference} concerns the Parisi measure $\mu_\beta$ itself, and
\cref{prop:lz-stability} concerns minimizing profiles. We need this generality,
since at the independent endpoint, the $n$-spin system carries the
profile of the full system.

\paragraph{Every minimizer.}
A minimizer need not be unique, and it may depend on $N$, $\rho$, and
$\lambda$, but not on the couplings. Since \eqref{eq:in-envelope} holds at
every minimizer, and since our bounds on $A$, $J$, and $V$ hold at every
minimizer with the same constants, we can integrate in $\rho$ without a
measurable choice of minimizers.

\paragraph{Atoms.}
Minimizing profiles may have atoms in $[0,1)$, and $\nu$ may charge
$0$, but $\nu(\{1\})=0$ by \cref{prop:fnd-contact}. Distribution functions are right-continuous, so that $a(q)$
includes an atom at $q$ while $(q,1]$ excludes it. We compute integrals
against $\nu$ through quantiles, and a quantile window may contain part
of an atom (\cref{sec:wi-setup}). We handle scalar and tree laws
through exact path likelihoods (\cref{lem:sc-likelihood,lem:tree-density})
and through the heat identity of \cref{lem:fnd-path-heat}, in which
profile measures are held fixed with their atoms. In
\cref{sec:lz,sec:lz-root}, we spread the free mass before using cutoffs.

\section{A map of the proof}\label{sec:overview-map}

\Cref{tab:overview-map} lists, for each section, its main results, its
main inputs, and the places where its results are used. The proof
of \cref{thm:main} in \cref{subsec:b7-general} begins with a table of
its inputs and of the parameters at which they are used. In terms of
these results, the proof has the following structure.
\begin{itemize}
\item \Cref{prop:fnd-mean-bias}, \cref{cor:fnd-defect}, and
\cref{lem:fnd-trial} reduce the theorem to a bound on the integral of
$\calD_N$, evaluated at any minimizer, and give the trial budget $\eta$
that defines the cutoff $\rho_0$.
\item On $(0,\rho_0]$, \cref{prop:sm-small} bounds the defect, using
\cref{prop:wi-vfs}, \cref{cor:lz-omega}, \eqref{eq:in-Lvar},
\cref{prop:th-j}, and \eqref{eq:fnd-moving}.
\item On $[\rho_0,1]$, \cref{cor:c7-defect} bounds the defect, using
\cref{prop:c7-closure}, \cref{cor:wi-cub}, and \cref{prop:th-j}. The
contraction of \cref{prop:c7-closure} combines
\cref{lem:cav-trace,lem:linear,lem:remainder,lem:cl-response,lem:cl-star,lem:c7-thermal},
and its stability hypothesis \eqref{eq:c7-stab} is verified by
\cref{prop:lz-stability}, which combines
\cref{lem:Jdom,lem:reference,lem:scalarcontinuity,prop:loc-isolation}.
\item The closure in \cref{subsec:b7-general} combines the two bounds.
\end{itemize}

\begin{table}[!tb]
\centering
\small
\begin{tabular}{@{}>{\raggedright\arraybackslash}p{0.06\textwidth}
 >{\raggedright\arraybackslash}p{0.31\textwidth}
 >{\raggedright\arraybackslash}p{0.3\textwidth}
 >{\raggedright\arraybackslash}p{0.25\textwidth}@{}}
\toprule
Sec. & Main results & Main inputs & Used in\\
\midrule
\ref{sec:fnd} & defect identity \eqref{eq:in-envelope}; contact and
 variance bounds, \cref{prop:fnd-contact,prop:fnd-variance}; trial
 budget, \cref{lem:fnd-trial}; reduction, \cref{prop:fnd-mean-bias}
 & \cite[Section~2, Lemma~7.2, Theorem~7.8]{AL},
 \cite{ChenConvexity} & all later sections\\
\ref{sec:gauss} & posterior bound for arbitrary profiles,
 \cref{lem:fnd-posterior}; heat identity, \cref{lem:fnd-path-heat}
 & \cite[Section~2]{AL},
 \cite[Lemma~2.3]{OpenAIFluctuationScale}
 & \cref{sec:sc,sec:c7,sec:c7-endpoint}\\
\ref{sec:wi} & variance on quantile intervals, \cref{prop:wi-vfs};
 cubic moment, \cref{cor:wi-cub}
 & \cref{prop:fnd-contact,prop:fnd-variance}
 & \cref{prop:sm-small,cor:c7-defect}\\
\ref{sec:th} & mean--covariance term, \cref{prop:th-j}
 & \cref{prop:fnd-contact,prop:fnd-variance}
 & \cref{prop:sm-small,cor:c7-defect}\\
\ref{sec:lz} & isolation, \cref{prop:loc-isolation}
 & \cite[Section~7.4, Appendix~A]{AL}, \cite[Theorem~1.1]{Lopatto},
 \cref{lem:fnd-trial}
 & \cref{sec:lz-root}, \cref{prop:lz-stability}\\
\ref{sec:lz-root} & root localization,
 \cref{prop:lz-quantile,cor:lz-omega}
 & \cref{prop:loc-inputs,lem:loc-smoothing,lem:loc-coercive}
 & \cref{prop:sm-small}\\
\ref{sec:sc} & response identity, \cref{lem:sc-response}; domination,
 \cref{lem:Jdom}; reference multiplier, \cref{lem:reference};
 stability, \cref{prop:lz-stability}
 & \cref{sec:gauss}, \cref{prop:loc-isolation},
 \cite{Lopatto,AuffingerChen2015Uniqueness}
 & \cref{lem:linear,prop:c7-closure}; proof of \cref{thm:main}\\
\ref{sec:c7} & tree calculus, \cref{lem:streaming,cor:comparison};
 remainder, \cref{lem:cav-threeproducts,lem:remainder}
 & \cref{sec:gauss}, scalar law of \cref{sec:sc}
 & \cref{sec:c7-endpoint,sec:c7-closure}\\
\ref{sec:c7-endpoint} & linear response, \cref{lem:linear}; errors,
 \cref{lem:cl-response}
 & \cref{sec:c7}, \eqref{eq:scalarJexplicit}, \cref{lem:fnd-path-heat}
 & \cref{lem:cav-trace,prop:c7-closure}\\
\ref{sec:c7-closure} & coefficient one, \cref{lem:cl-star}; trace
 identity, \cref{lem:cav-trace}; contraction, \cref{prop:c7-closure};
 defect, \cref{cor:c7-defect}
 & \cref{sec:c7,sec:c7-endpoint}, \cref{lem:Jdom},
 \cref{cor:wi-cub,prop:th-j}
 & proof of \cref{thm:main}\\
\ref{sec:b7} & small $\rho$, \cref{prop:sm-small}; \cref{thm:main}
 & \cref{sec:fnd,sec:wi,sec:th,sec:lz-root,sec:sc,sec:c7-closure}
 & ---\\
\bottomrule
\end{tabular}
\caption{The sections of the proof, their main results and inputs, and
where the results are used.}
\label{tab:overview-map}
\end{table}
